\chapter{Varsayımlar ve Parametrik Olmayan Alternatifler}
\index{parametrik olmayan yöntem}
\index{varsayım kontrolü}

\begin{hedefler}
Bu bölümün sonunda okuyucu:
\begin{itemize}
  \item varsayım kontrolünü tek bir normallik testine indirgemeden araştırma tasarımıyla birlikte değerlendirebilecek,
  \item parametrik olmayan yöntemlerin hedefini ve varsayımlarını açıklayabilecek,
  \item bağımsız iki grup için Mann--Whitney testini uygulayıp doğru yorumlayabilecek,
  \item eşleştirilmiş veride işaret testi ile Wilcoxon işaretli sıra testini ayırabilecek,
  \item üç veya daha fazla bağımsız grup için Kruskal--Wallis testini ve düzeltilmiş takip analizlerini kullanabilecek,
  \item kesin ve asimptotik $p$-değerleri ile bağlı sıraların rolünü tanıyabilecek,
  \item sıra tabanlı sonuçları medyan farkı gibi aşırı yorumlamadan uygun etki büyüklüğüyle raporlayabilecek,
  \item dönüşüm, sağlam yöntem, permütasyon ve bootstrap seçeneklerini araştırma hedefine göre karşılaştırabilecektir.
\end{itemize}
\end{hedefler}

\begin{hazirlik}
``Veri normal değilse her zaman parametrik olmayan test kullanılır'' önermesini
örneklem hacmi, aykırı değer, bağımsızlık ve araştırmanın ortalama mı yoksa
sıralama mı hedeflediği açısından değerlendirin.
\end{hazirlik}

\section{PDR vakası: Yardım aramaya hazır oluş puanları}

Bir psikolojik danışma merkezi, iki bilgilendirme yönteminin öğrencilerin
yardım aramaya hazır oluş düzeyiyle ilişkisini incelemektedir. Ölçüm 1 ile 10
arasında sıralı bir derecelendirmedir. Yüz yüze bilgilendirme grubunun puanları
\[
7,8,6,9,7,8,10,6,9,8
\]
ve çevrim içi bilgilendirme grubunun puanları
\[
4,6,5,7,5,6,4,8,5,6
\]
olarak gözlenmiştir.

Ölçek yalnız sıralı yorumlanabiliyorsa ortalama farkı yerine öğrencilerin
göreli sıralarını karşılaştırmak daha uygun olabilir. Bununla birlikte
``normal değil, o hâlde Mann--Whitney'' biçimindeki otomatik karar yeterli
değildir. Önce şu sorular sorulmalıdır:

\begin{enumerate}
  \item Gruplar gerçekten bağımsız mıdır?
  \item Puanların büyüklükleri mi, sıraları mı bilimsel olarak anlamlıdır?
  \item Dağılımların biçimleri benzer midir?
  \item Çok sayıda eşit puan, yani bağlı sıra var mıdır?
  \item Araştırma hedefi dağılım, konum, medyan veya üstünlük olasılığından hangisidir?
\end{enumerate}

Bu bölümde parametrik olmayan testler, ``kolay yedek yöntem'' olarak değil,
belirli hedef büyüklüklere ve varsayımlara sahip istatistiksel araçlar olarak
ele alınacaktır.

\begin{researchbox}
Hazır oluş puanı sıralı olsa bile kategoriler arası psikolojik mesafe eşit
olmayabilir. 3'ten 4'e geçiş ile 8'den 9'a geçiş aynı değişimi temsil etmiyorsa
ortalama farkının yorumu tartışmalıdır. Yöntem seçimi yazılımın değişkeni
``scale'' olarak işaretlemesine değil, ölçme kuramına dayanmalıdır.
\end{researchbox}

\section{Yöntem adı değil araştırma hedefi}
\index{hedef büyüklük}
\index{ölçme düzeyi}

Parametrik ve parametrik olmayan yöntemler aynı sorunun her koşulda birbiriyle
değiştirilebilir iki sürümü değildir. Hedefledikleri nicelikler farklı olabilir
\cite{hollander2013,lehmann2006,wasserman2006}.

\begin{table}[htbp]
\centering
\caption{Araştırma yapısı ve yaygın analiz seçenekleri.}
\begin{tabular}{@{}p{0.27\textwidth}p{0.25\textwidth}p{0.29\textwidth}@{}}
\toprule
Araştırma yapısı & Ortalama temelli yöntem & Sıra veya işaret temelli seçenek \\
\midrule
Tek ya da eşleştirilmiş fark & Tek veya eşleştirilmiş $t$ testi & İşaret veya Wilcoxon testi \\
İki bağımsız grup & Welch $t$ testi & Mann--Whitney testi \\
Üçten fazla bağımsız grup & Welch/klasik ANOVA & Kruskal--Wallis testi \\
İki değişkenin ilişkisi & Pearson korelasyonu & Spearman korelasyonu \\
\bottomrule
\end{tabular}
\end{table}

Araştırma sorusu ``ortalama kaç puan değişti?'' ise sıra testi doğrudan aynı
hedefi sınamayabilir. Ağır çarpıklık durumunda ortalama için sağlam standart
hata, permütasyon testi veya bootstrap güven aralığı hedefi daha iyi
koruyabilir. Sıralama davranışı araştırma açısından anlamlıysa sıra tabanlı
yöntem doğal seçimdir.

\begin{mistakebox}
Parametrik olmayan test ``varsayımsız test'' demek değildir. Bağımsızlık,
eşleştirme yapısı, ölçme düzeyi, dağılım biçimi ve bağlı sıralar yine önemlidir.
Yalnız normal dağılım varsayımı farklılaştırılmış olur.
\end{mistakebox}

\section{Varsayım kontrolünde dengeli yaklaşım}
\index{varsayım kontrolü}
\index{Shapiro--Wilk testi}
\index{Q--Q grafiği}

Bir yöntemin uygunluğu şu sırayla değerlendirilmelidir:

\begin{enumerate}
  \item Araştırma tasarımı ve analiz birimi doğru mu?
  \item Hedef parametre veya karşılaştırma açık mı?
  \item Ölçme düzeyi yapılan işlemleri destekliyor mu?
  \item Grafiklerde çarpıklık, ağır kuyruk, kümelenme ve aykırı değer var mı?
  \item Örneklem hacmi ve grup dengesi ne düzeyde?
  \item Seçilen yöntemin sonuç yorumu araştırma sorusuyla uyumlu mu?
\end{enumerate}

Shapiro--Wilk gibi normallik testleri bu değerlendirmenin yalnız küçük bir
parçasıdır. Büyük örneklemde önemsiz sapmalar düşük $p$-değeri, küçük örneklemde
ciddi sapmalar yüksek $p$-değeri üretebilir. Histogram ve Q--Q grafiği de tek
başına kesin karar vermez; örneklem büyüklüğü ve yöntemin sağlamlığıyla birlikte
okunur.

\begin{assumptionbox}
Bağımsızlık ihlali dağılım dönüşümüyle veya sıra testine geçerek düzelmez.
Aynı öğrencinin iki ölçümü, öğrencilerin sınıflar içinde kümelenmesi veya
eşleştirilmiş çiftler tasarıma uygun yöntem gerektirir.
\end{assumptionbox}

\section{Parametrik yöntemler ne zaman sağlamdır?}
\index{sağlamlık}
\index{Welch t testi@Welch $t$ testi}

$t$ testi ve ANOVA, özellikle gruplar dengeli ve örneklemler yeterince büyükse
ılımlı normallik sapmalarına dayanıklı olabilir. Buna karşılık küçük örneklem,
güçlü çarpıklık, ağır kuyruk ve etkili aykırı değer birlikte bulunduğunda
ortalama ile standart hata kararsızlaşabilir.

Varyans eşitsizliği normal olmama ile aynı sorun değildir. Bağımsız iki grupta
Welch testi, çok grupta Welch ANOVA değişen varyansa doğrudan çözüm sunar.
Yalnız Levene testi anlamlı diye sıra testine geçmek hedefi gereksiz biçimde
değiştirebilir.

\begin{sezgibox}
Varsayım kontrolünün amacı ``parametrik testi geçmek'' değildir. Amaç,
hesaplanan tahminin ve belirsizliğin araştırma sorusunu güvenilir biçimde
temsil edip etmediğini anlamaktır.
\end{sezgibox}

\section{Sıralama mantığı ve bağlı sıralar}
\index{sıra}
\index{bağlı sıra}

Sıra tabanlı yöntemlerde bütün gözlemler küçükten büyüğe dizilir. En küçük
değere 1, sonraki değere 2 sırası verilir. Eşit değerler aynı sıraları
paylaşır ve hak edecekleri sıraların ortalamasını alır.

Örneğin sıralı değerler $2,4,4,7$ ise sıralar
\[
1,\ 2{,}5,\ 2{,}5,\ 4
\]
olur. Çok sayıda bağlı sıra test istatistiğinin varyansını etkiler; yazılımlar
bağ düzeltmesi uygular. El hesabındaki basit formüller bağlı sıra yokmuş gibi
kullanılmamalıdır.

Sıralama uç değerlerin büyüklük etkisini azaltır. 20 yerine 200 gözlenmesi en
büyük sıra konumunu değiştirmiyorsa sıra testi aynı veya benzer sonuç verebilir.
Bu sağlamlık, değerler arasındaki mesafe bilgisinin kullanılmaması karşılığında
elde edilir.

\section{Mann--Whitney testi}
\index{Mann--Whitney testi}
\index{üstünlük olasılığı}
\index{sıra-biserial korelasyon}

Mann--Whitney testi iki bağımsız grubun gözlemlerini birlikte sıralar ve sıra
toplamlarını karşılaştırır. Sıfır hipotezi genel biçimde iki grubun dağılımının
aynı olmasıdır. Test istatistiği, gruplar arası bütün gözlem çiftlerindeki
üstünlük sayısıyla ilişkilidir.

Birinci gruptaki $X$, ikinci gruptaki $Y$ gözlemi için
\[
\theta=P(X>Y)+\frac12P(X=Y)
\]
olasılığı düşünülebilir. Sıfır altında $\theta=0{,}50$'dir. Bir grubun rastgele
seçilen üyesinin diğer gruptakinden daha yüksek puan alma eğilimi testin doğal
yorumlarından biridir.

\subsection{Adım adım bağımsız grup örneği}

Hazır oluş verisinde medyan ve çeyrekler şöyledir:

\begin{table}[htbp]
\centering
\caption{İki bilgilendirme grubunun betimsel özetleri.}
\begin{tabular}{@{}lrrr@{}}
\toprule
Grup & $n$ & Medyan & $Q_1$--$Q_3$ \\
\midrule
Yüz yüze & 10 & 8,0 & 7,0--8,75 \\
Çevrim içi & 10 & 5,5 & 5,0--6,0 \\
\bottomrule
\end{tabular}
\end{table}

Yazılımın birinci grup için verdiği istatistik $U=88{,}5$ ve iki yönlü
$p=0{,}004$'tür. $U$'nun hangi grup için ve büyük ya da küçük sürüm olarak
raporlandığı yazılıma göre değişebilir; $p$-değeri ve grup sıralarıyla birlikte
okunmalıdır.

Üstünlük olasılığı tahmini
\[
\hat\theta=\frac{U}{n_1n_2}
=\frac{88{,}5}{100}=0{,}885
\]
olur. Bağların yarısı paylaştırıldığında yüz yüze gruptan rastgele seçilen bir
öğrencinin çevrim içi gruptakinden daha yüksek hazır oluş puanı alma olasılığı
yaklaşık yüzde 88,5'tir.

Sıra-biserial etki büyüklüğü
\[
r_{rb}=2\hat\theta-1=0{,}77
\]
olarak hesaplanabilir. İşaret grup sırasına bağlıdır; kullanılan yön açıkça
belirtilmelidir.

\begin{mistakebox}
Mann--Whitney testi her durumda ``iki medyan eşit mi?'' testi değildir.
Dağılım biçimleri ve yayılımları farklıysa anlamlı sonuç konum, yayılım veya
genel dağılım farkından kaynaklanabilir. Medyan farkı yorumu benzer dağılım
biçimi gibi ek koşullar gerektirir.
\end{mistakebox}

\section{Mann--Whitney için koşullar ve grafikler}

Mann--Whitney testinde:

\begin{itemize}
  \item iki grup birbirinden bağımsızdır,
  \item gözlemler grup içinde bağımsızdır,
  \item sonuç en az sıralı ölçektedir,
  \item üstünlük yorumu için gözlemler ortak bir anlamlı sıralamada karşılaştırılabilir,
  \item yalnız konum veya medyan farkı yorumu yapılacaksa dağılım biçimleri yeterince benzerdir.
\end{itemize}

Yan yana kutu, nokta veya violin grafikleri yalnız medyanı değil yayılımı,
çarpıklığı ve bağlı puanları da göstermelidir. Aynı medyana sahip iki grup çok
farklı dağılımlar nedeniyle anlamlı Mann--Whitney sonucu verebilir.

\section{İşaret testi}
\index{işaret testi}
\index{binom testi}

Tek örneklem veya eşleştirilmiş veride işaret testi yalnız farkların yönünü
kullanır. Sıfır farklar analizden çıkarıldıktan sonra pozitif fark sayısı,
$H_0$ altında başarı olasılığı 0,5 olan binom dağılımıyla karşılaştırılır.

On iki öğrencinin son-test eksi ön-test farkları
\[
-4,-1,-4,-2,-3,-1,-3,-4,1,-5,-2,-4
\]
olsun. On bir negatif ve bir pozitif fark vardır. İki yönlü kesin işaret testi
$p\approx0{,}006$ verir.

İşaret testi fark büyüklüklerini kullanmadığı için aykırı büyüklüklere çok
dayanıklıdır; ancak bilgi kaybı nedeniyle Wilcoxon testinden daha düşük güce
sahip olabilir. Doğal hedef eşleştirilmiş farkların medyanı veya pozitif fark
olasılığıdır.

\section{Wilcoxon işaretli sıra testi}
\index{Wilcoxon işaretli sıra testi}
\index{fark dağılımı!simetri}

Wilcoxon testi sıfır olmayan farkların mutlak değerlerini sıralar, sonra
sıralara fark işaretlerini geri verir. Pozitif ve negatif sıra toplamları
karşılaştırılır. Sıfır hipotezi, fark dağılımının sıfır çevresinde simetrik
olmasıdır.

Örnekte fark medyanı $-3$'tür. Pozitif sıra toplamı 2, negatif sıra toplamı
76 ve yazılımın küçük toplam olarak verdiği istatistik $W=2$'dir. İki yönlü
kesin sonuç
\[
p=0{,}002
\]
olarak bulunur. Son-test puanları çoğunlukla ön-testten düşüktür.

Eşleştirilmiş sıra-biserial etki büyüklüğü
\[
r_{rb}=\frac{W_+-W_-}{W_++W_-}
=\frac{2-76}{78}\approx-0{,}949
\]
olur. Negatif yön, farkın ``sonra eksi önce'' tanımına bağlıdır.

\begin{mistakebox}
Wilcoxon testi farkların normal olmasını gerektirmez; ancak konum değişimi
yorumu için fark dağılımının yaklaşık simetrik olması önemlidir. Güçlü
asimetrik farklarda işaret testi veya başka sağlam yöntem daha uygun olabilir.
\end{mistakebox}

\section{Sıfır farklar ve yazılım seçenekleri}
\index{sıfır fark}
\index{Pratt yöntemi}

Eşleştirilmiş ölçümler aynıysa fark sıfırdır. Klasik Wilcoxon uygulaması sıfır
farkları çıkarabilir; Pratt gibi yaklaşımlar sıralama sürecinde farklı biçimde
ele alır. Yazılımın varsayılanı sonuçları, özellikle küçük örneklemde,
değiştirebilir.

Rapor geçerli çift sayısını, sıfır fark sayısını, kesin veya asimptotik
$p$-değerini ve kullanılan yöntemi belirtmelidir. Çok sayıda sıfır fark,
ölçeğin değişime duyarlılığı veya taban--tavan etkisi hakkında da bilgi verir.

\section{Kruskal--Wallis testi}
\index{Kruskal--Wallis testi}
\index{sıra ortalaması}

Kruskal--Wallis testi üç veya daha fazla bağımsız grubun ortak sıra
dağılımını karşılaştırır. Bütün gözlemler birlikte sıralanır ve grup sıra
ortalamalarının genel sıra ortalamasından uzaklığı ölçülür.

Bağ düzeltmesi bulunmayan temel istatistik
\[
H=\frac{12}{N(N+1)}
\sum_{j=1}^k\frac{R_j^2}{n_j}-3(N+1)
\]
biçimindedir. Burada $R_j$ grup sıra toplamıdır. Büyük örneklemde $H$ yaklaşık
$k-1$ serbestlik dereceli ki-kare dağılımıyla karşılaştırılır.

\subsection{Üç program örneği}

Üç bağımsız program grubunun 1--10 arası hazır oluş puanlarında medyanlar
sırasıyla 8, 6 ve 4,5 olsun. Test sonucu
\[
H(2)=21{,}91,
\qquad p<0{,}001
\]
olarak bulunur. En az bir grubun sıra dağılımı farklıdır; hangi grupların
farklı olduğu henüz bilinmez.

Kruskal--Wallis için epsilon-kare etki büyüklüğünün bir sürümü
\[
\varepsilon^2=\frac{H-k+1}{N-k}
\]
olup bu örnekte
\[
\varepsilon^2
=\frac{21{,}91-3+1}{30-3}
\approx0{,}737
\]
bulunur. Ölçü örneklemde güçlü sıra ayrışmasına işaret eder; evrensel eşik
yerine bağlam ve güven aralığı tercih edilmelidir.

\section{Genel testten sonra ikili karşılaştırmalar}
\index{Dunn testi}
\index{Holm düzeltmesi}

Anlamlı Kruskal--Wallis sonucu farkın yerini göstermez. Dunn testi veya ikili
Mann--Whitney testleri çoklu karşılaştırma düzeltmesiyle kullanılabilir.
Örnekte ham ikili $p$-değerleri yaklaşık 0,0010, 0,0002 ve 0,0029'dur. Holm
düzeltmesinden sonra üç fark da yüzde 5 aile düzeyinde belirgin kalmaktadır.

Yalnız genel test anlamlı diye düzeltmesiz bütün ikilileri sınamak yanlış
pozitif oranını yükseltir. Karşılaştırmalar veri görülmeden önce planlanmışsa
hedefli kontrast yaklaşımı daha verimli olabilir.

\begin{mistakebox}
Bir çiftin anlamlı, diğerinin anlamsız olması bu iki etkinin birbirinden
anlamlı biçimde farklı olduğunu kanıtlamaz. Her hedef farkın doğrudan tahmini
ve düzeltilmiş karşılaştırması incelenmelidir.
\end{mistakebox}

\section{Kesin ve asimptotik p-değerleri}
\index{kesin p-değeri}
\index{asimptotik p-değeri}
\index{Monte Carlo yöntemi}

Küçük örneklemde test istatistiğinin olası bütün değerleri sayılarak kesin
$p$-değeri hesaplanabilir. Büyük örneklemde normal veya ki-kare yaklaşımıyla
asimptotik sonuç elde edilir. Bağlı sıralar ve sıfır farklar kesin dağılımı
karmaşıklaştırabilir; yazılım koşullu kesin, Monte Carlo veya düzeltilmiş
asimptotik yöntem kullanabilir.

Kesin ve asimptotik sonuçlar küçük farklılık gösterebilir. Araştırmacı yalnız
daha düşük olanı seçmemeli; örneklem yapısına uygun yöntemi önceden belirleyip
raporlamalıdır.

\section{Etki büyüklüğü ve güven aralığı}
\index{epsilon-kare@$\varepsilon^2$}
\index{Hodges--Lehmann tahmini}

Parametrik olmayan testte $p$-değeri örneklem büyüklüğünden etkilenir. Etki
büyüklüğü ilişkinin derecesini görünür kılar. Yaygın seçenekler:

\begin{itemize}
  \item Mann--Whitney için üstünlük olasılığı, Cliff deltası veya sıra-biserial korelasyon,
  \item Wilcoxon için eşleştirilmiş sıra-biserial korelasyon,
  \item Kruskal--Wallis için epsilon-kare,
  \item medyan ya da konum farkı için Hodges--Lehmann tahmini ve güven aralığı.
\end{itemize}

Sıra-biserial korelasyon Pearson korelasyonuyla aynı bilimsel nicelik değildir;
yalnız ölçeği $-1$ ile 1 arasındadır. Etki ölçüsünün adı ve yön tanımı açıkça
verilmelidir.

\section{Parametrik ve parametrik olmayan sonuçlar farklıysa}

İki yöntem farklı hedefleri ve veri özelliklerini ağırlıklandırdığı için
farklı sonuç verebilir. Olası nedenler:

\begin{itemize}
  \item aykırı değerlerin ortalamayı güçlü etkilemesi,
  \item grupların farklı yayılım veya biçime sahip olması,
  \item sıra testinin değerler arasındaki mesafeyi kullanmaması,
  \item küçük örneklem ve düşük güç,
  \item bağlı sıralar veya tavan--taban etkileri.
\end{itemize}

Araştırmacı istenen sonucu veren testi seçmemelidir. Ana hedef ve yöntem
önceden belirlenmeli; alternatif analiz duyarlılık incelemesi olarak ve neden
farklılaştığı açıklanarak sunulmalıdır.

\section{Dönüşüm, sağlam yöntem, permütasyon ve bootstrap}
\index{dönüşüm}
\index{sağlam yöntem}
\index{permütasyon testi}
\index{bootstrap}

Sıra testine geçmek tek seçenek değildir.

\begin{description}
  \item[Dönüşüm] Sağa çarpık sürelerde logaritma gibi dönüşüm model yapısını iyileştirebilir; yorum dönüştürülmüş ölçektedir.
  \item[Sağlam yöntem] Kırpılmış ortalama ve sağlam standart hatalar aykırı değer etkisini azaltabilir.
  \item[Permütasyon] Grup etiketlerinin değiştirilebilirliği altında hedef istatistiğin sıfır dağılımını üretir.
  \item[Bootstrap] Gözlenen örneklemden yeniden örnekleyerek tahminin belirsizliğini değerlendirir.
\end{description}

Permütasyon ve bootstrap ``varsayımsız'' değildir. Değiştirilebilirlik,
bağımsızlık ve örneklemin hedef evreni temsil etmesi gibi koşullar sürer.
Kötü örnekleme tasarımı yeniden örneklemeyle düzelmez
\cite{efron1993,wasserman2006}.

\begin{researchbox}
PDR araştırmasında yöntem seçimi yalnız istatistiksel uygunluk değildir.
Medyan değişim, öğrencilerin çoğunun yönü, klinik eşik aşımı veya ortalama
puan değişimi farklı hizmet sorularını yanıtlar. Önce kararın hangi niceliğe
dayanacağı belirlenmelidir.
\end{researchbox}

\section{Aykırı değerleri silmek yerine anlamak}

Parametrik olmayan teste geçmek aykırı gözlemi araştırma sorumluluğunu ortadan
kaldırmaz. Uç değer veri girişi hatası, farklı bir alt grup veya gerçek ve
önemli bir danışan durumu olabilir. İnceleme sırası:

\begin{enumerate}
  \item özgün kayıt ve ölçü birimi doğrulanır,
  \item katılımcının uygunluk ölçütleri kontrol edilir,
  \item sonuç ortalama, medyan ve sıra yöntemleriyle karşılaştırılır,
  \item silme kararı sonuçtan bağımsız gerekçelendirilir,
  \item duyarlılık sonuçları şeffaf biçimde raporlanır.
\end{enumerate}

\section{Python ile parametrik olmayan testler}

\begin{lstlisting}
import numpy as np
from scipy import stats

# Iki bagimsiz grup
yuz_yuze = np.array([7,8,6,9,7,8,10,6,9,8])
cevrim_ici = np.array([4,6,5,7,5,6,4,8,5,6])
mw = stats.mannwhitneyu(
    yuz_yuze, cevrim_ici, alternative="two-sided"
)
ustunluk = mw.statistic / (len(yuz_yuze) * len(cevrim_ici))
rank_biserial = 2 * ustunluk - 1

# Eslestirilmis farklar
fark = np.array([-4,-1,-4,-2,-3,-1,-3,-4,1,-5,-2,-4])
wilcoxon = stats.wilcoxon(fark, alternative="two-sided")
isaret = stats.binomtest(
    k=(fark > 0).sum(), n=(fark != 0).sum(), p=0.5
)

# Uc bagimsiz grup
A = np.array([8,9,7,8,10,9,8,7,9,8])
B = np.array([6,7,5,8,6,7,6,5,7,6])
C = np.array([4,5,6,3,5,4,6,4,5,3])
kw = stats.kruskal(A, B, C)
epsilon2 = (kw.statistic - 3 + 1) / (30 - 3)

print("Mann-Whitney:", mw)
print("Ustunluk ve sıra-biserial:", ustunluk, rank_biserial)
print("Wilcoxon:", wilcoxon)
print("Isaret testi:", isaret)
print("Kruskal-Wallis:", kw)
print("Epsilon-kare:", epsilon2)
\end{lstlisting}

Yazılım sürümüne göre kesin ve asimptotik yöntemin varsayılanı değişebilir.
Bağlı sıralar, sıfır farklar ve alternatif hipotez seçeneği kontrol edilmelidir.
İkili Kruskal--Wallis takipleri için düzeltilmiş $p$-değerleri ayrıca
hesaplanmalıdır.

\section{SPSS ile uygulama ve çıktı okuma}

\begin{enumerate}
  \item Bağımsız gruplar için \textbf{Analyze $>$ Nonparametric Tests $>$ Independent Samples} yolunu izleyin.
  \item Eşleştirilmiş ölçümler için \textbf{Related Samples} bölümünde Wilcoxon veya işaret testini seçin.
  \item Üç veya daha fazla bağımsız grup için Kruskal--Wallis testini ve gerekiyorsa ikili karşılaştırmaları etkinleştirin.
  \item Değişkenlerin grup kodlarını, ölçüm çiftlerini ve kategori sırasını doğrulayın.
  \item Exact veya Monte Carlo seçeneği kullanıldıysa yöntemi ve rassal ayarları kaydedin.
  \item Test istatistiği yanında grup medyanları, IQR, sıra ortalamaları ve etki büyüklüğünü raporlayın.
\end{enumerate}

\begin{outputguidebox}
Çıktıyı şu sırayla okuyun: (1) testin veri yapısına uygunluğu, (2) geçerli grup
veya çift sayıları, (3) medyan, IQR ve grafikler, (4) sıra toplamları veya sıra
ortalamaları, (5) test istatistiği, kesin/asimptotik yöntem ve $p$-değeri,
(6) uygun etki büyüklüğü, (7) genel testten sonra düzeltilmiş karşılaştırmalar.
Yalnız ``Sig.'' hücresini raporlamayın.
\end{outputguidebox}

\section{Üç raporlama örneği}

\subsection*{Mann--Whitney}

``Yardım aramaya hazır oluş puanı yüz yüze grupta daha yüksektir
($Md=8{,}0$, $IQR=7{,}0$--$8{,}75$) ve çevrim içi grupta daha düşüktür
($Md=5{,}5$, $IQR=5{,}0$--$6{,}0$), $U=88{,}5$, $p=0{,}004$.
Üstünlük olasılığı 0,885 ve sıra-biserial etki $r_{rb}=0{,}77$'dir.''

\subsection*{Wilcoxon}

``Son-test eksi ön-test farklarının medyanı $-3$ puandır. Son-test puanları
çoğunlukla daha düşüktür, $W=2$, kesin $p=0{,}002$,
$r_{rb}=-0{,}949$. Fark yönü sonra eksi önce olarak tanımlanmıştır.''

\subsection*{Kruskal--Wallis}

``Üç programın hazır oluş sıra dağılımları farklıdır,
$H(2)=21{,}91$, $p<0{,}001$, $\varepsilon^2=0{,}737$. Holm düzeltmeli ikili
karşılaştırmalar bütün grup çiftlerinde fark göstermiştir. Medyan ve IQR
değerleri her grup için ayrıca sunulmuştur.''

Bu ifadeler nedensellik iddiası taşımaz. Gruplara rastgele atama yapılmadıysa
gözlenen sıra farkları başlangıç özellikleriyle açıklanabilir.

\section{Bir parametrik olmayan analizi değerlendirirken}

\begin{enumerate}
  \item Test araştırma hedefi ve veri yapısıyla mı, yalnız normallik $p$-değeriyle mi seçilmiş?
  \item Bağımsız veya eşleştirilmiş yapı doğru tanımlanmış mı?
  \item Grafikler dağılım biçimi, yayılım ve aykırı değerleri gösteriyor mu?
  \item Mann--Whitney sonucu gerekçesiz biçimde medyan farkı olarak mı yorumlanmış?
  \item Wilcoxon için farkların simetrisi, sıfırlar ve bağlı sıralar incelenmiş mi?
  \item Kruskal--Wallis sonrası karşılaştırmalar çoklu test için düzeltilmiş mi?
  \item Kesin veya asimptotik yöntemin hangisi olduğu belirtilmiş mi?
  \item Etki büyüklüğü, medyan ve IQR $p$-değeriyle birlikte verilmiş mi?
  \item Alternatif analizler istenen sonucu seçmek için mi kullanılmış?
\end{enumerate}

\section{Literatür veri setiyle IBM SPSS laboratuvarı}
\index{PlantGrowth veri seti}
\index{SPSS!Kruskal--Wallis testi}
\index{Dunn--Bonferroni karşılaştırması}
\index{epsilon-kare@$\varepsilon^2$!SPSS uygulaması}

Bu laboratuvarda R'nin \texttt{datasets} paketinde bulunan klasik
\emph{PlantGrowth} verisi kullanılacaktır \cite{dobson1983}. Veri seti,
kontrol koşulu ile iki farklı uygulama koşulunda yetiştirilen bitkilerin kuru
ağırlıklarını içerir. Her grupta 10 olmak üzere toplam 30 bağımsız gözlem
vardır. Özgün değişken adları \texttt{weight} ve \texttt{group} korunmuştur.

Uygulamanın araştırma sorusu ``Üç koşuldaki kuru ağırlıkların sıra dağılımları
aynı mıdır?'' biçimindedir. Hipotezler
\[
H_0:\text{üç grubun kuru ağırlık dağılımları aynıdır},
\]
\[
H_1:\text{en az bir grubun dağılımı farklıdır}
\]
olarak kurulacaktır. Grupların yalnız konum bakımından karşılaştırıldığı bir
yorum yapılacaksa dağılım biçimlerinin de yeterince benzer olduğu grafiklerle
değerlendirilmelidir.

\begin{researchbox}
Bu veri PDR alanından değil, istatistik öğretiminde kullanılan deneysel bir
örnektir. Buradaki amaç bağlamlar arasında bilimsel genelleme yapmak değil;
bağımsız üç grubun sıra temelli karşılaştırmasını gerçek ve yayımlanmış küçük
bir veri seti üzerinde yeniden üretmektir. Bitki kuru ağırlığı yerine PDR
araştırmasında sıralı hazır oluş, memnuniyet veya belirgin biçimde çarpık bir
süre değişkeni aynı analiz yapısına sahip olabilir.
\end{researchbox}

\subsection{Verinin SPSS'e girilmesi}

Aşağıdaki syntax 30 gözlemi uzun veri biçiminde oluşturur. Her satır tek bir
bitkiyi, \texttt{group} değişkeni ise kontrol, tedavi 1 veya tedavi 2 koşulunu
gösterir.

\begin{lstlisting}[language={}]
DATA LIST LIST /weight (F4.2) group (F1.0).
BEGIN DATA
4.17 1
5.58 1
5.18 1
6.11 1
4.50 1
4.61 1
5.17 1
4.53 1
5.33 1
5.14 1
4.81 2
4.17 2
4.41 2
3.59 2
5.87 2
3.83 2
6.03 2
4.89 2
4.32 2
4.69 2
6.31 3
5.12 3
5.54 3
5.50 3
5.37 3
5.29 3
4.92 3
6.15 3
5.80 3
5.26 3
END DATA.

DATASET NAME PlantGrowth WINDOW=FRONT.
VARIABLE LABELS
 weight 'Bitkinin kuru ağırlığı'
 group  'Deney koşulu'.
VALUE LABELS group
 1 'Kontrol' 2 'Tedavi 1' 3 'Tedavi 2'.
VARIABLE LEVEL weight (SCALE) group (NOMINAL).
EXECUTE.
\end{lstlisting}

\subsection{Betimsel inceleme ve genel test}

Önce grup büyüklükleri, medyanlar ve kutu grafikleri incelenir. Ardından
Kruskal--Wallis testi yürütülür. İlk komut eski diyaloglarla uyumlu genel
testi; ikinci komut yeni parametrik olmayan testler yordamında Bonferroni
düzeltmeli Dunn ikili karşılaştırmalarını üretir.

\begin{lstlisting}[language={}]
EXAMINE VARIABLES=weight BY group
 /PLOT BOXPLOT
 /STATISTICS DESCRIPTIVES
 /PERCENTILES(25,50,75)
 /MISSING LISTWISE
 /NOTOTAL.

NPAR TESTS
 /K-W=weight BY group(1 3)
 /MISSING ANALYSIS.

NPTESTS
 /INDEPENDENT TEST (weight) GROUP (group)
  KRUSKAL_WALLIS(COMPARE=PAIRWISE)
 /MISSING SCOPE=ANALYSIS USERMISSING=EXCLUDE
 /CRITERIA ALPHA=0.05 CILEVEL=95.
\end{lstlisting}

Menüyle çalışmak için \textbf{Analyze $>$ Nonparametric Tests $>$ Independent
Samples} yolu izlenir. \textbf{Fields} sekmesinde \texttt{weight} test,
\texttt{group} grup değişkeni olarak tanımlanır. \textbf{Settings} altında
testler özelleştirilerek Kruskal--Wallis ve bütün ikili karşılaştırmalar
seçilir. SPSS sürümüne göre karşılaştırma ayarlarının adı ve çıktının görünümü
değişebileceğinden üretilen syntax analiz dosyasıyla birlikte saklanmalıdır.

\subsection{Grup özetleri ve sıra ortalamaları}

\begin{table}[htbp]
\centering
\caption{PlantGrowth gruplarının betimsel ve sıra temelli özetleri.}
\begin{tabular}{@{}lrrrr@{}}
\toprule
Grup & $n$ & Medyan & Ortalama & Sıra ortalaması \\
\midrule
Kontrol & 10 & 5,155 & 5,032 & 14,75 \\
Tedavi 1 & 10 & 4,550 & 4,661 & 10,35 \\
Tedavi 2 & 10 & 5,435 & 5,526 & 21,40 \\
\bottomrule
\end{tabular}
\end{table}

Örneklemde en düşük medyan ve sıra ortalaması tedavi 1, en yüksek değerler
tedavi 2 grubundadır. Kontrol grubu ikisinin arasında yer alır. Bununla
birlikte bu sıralama tek başına evren dağılımlarının farklı olduğunu
göstermez; grup içi örtüşme genel testte dikkate alınmalıdır.

Kutu grafiklerinde yalnız merkez değil, yayılım ve dağılım biçimi de
karşılaştırılır. Örneklem küçük olduğundan grafiklerdeki benzerlik veya
farklılık kesin bir varsayım kanıtı olarak değil, test sonucunun hangi tür
dağılım ayrışmasını temsil ettiğini anlamaya yardımcı bilgi olarak kullanılır.

\subsection{Kruskal--Wallis çıktısının analizi}

SPSS çıktısında Kruskal--Wallis istatistiği
\[
H(2)=7{,}988,
\qquad p=0{,}018
\]
olarak bulunur. Yüzde 5 anlamlılık düzeyinde üç grubun ortak sıra dağılımının
aynı olduğu hipoteze karşı kanıt vardır. Sonuç bütün grup çiftlerinin farklı
olduğunu göstermez; farkın yeri düzeltilmiş ikili karşılaştırmalarla
incelenmelidir.

Bu uygulamada epsilon-kare
\[
\varepsilon^2
=\frac{H-k+1}{N-k}
=\frac{7{,}988-3+1}{30-3}
\approx0{,}222
\]
olarak hesaplanır. Bu değer örneklemde gruplar arasında dikkate değer bir sıra
ayrışması bulunduğunu gösterir; evrensel ``küçük--orta--büyük'' eşikleriyle
mekanik biçimde sınıflandırılmamalıdır.

\subsection{Dunn--Bonferroni ikili karşılaştırmaları}

\begin{table}[htbp]
\centering
\caption{Kruskal--Wallis sonrası Dunn ikili karşılaştırmaları.}
\begin{tabular}{@{}lrrr@{}}
\toprule
Karşılaştırma & Sıra ortalaması farkı & Ham $p$ & Düzeltilmiş $p$ \\
\midrule
Kontrol--Tedavi 1 & 4,40 & 0,264 & 0,791 \\
Kontrol--Tedavi 2 & $-6{,}65$ & 0,091 & 0,273 \\
Tedavi 1--Tedavi 2 & $-11{,}05$ & 0,005 & 0,015 \\
\bottomrule
\end{tabular}
\end{table}

Bonferroni düzeltmesinden sonra yalnız tedavi 1 ile tedavi 2 arasındaki
karşılaştırma yüzde 5 düzeyinde belirgindir. Negatif sıra farkı, tablo yönü
``tedavi 1 eksi tedavi 2'' olduğu için tedavi 1 grubunun daha düşük sıralara
sahip olduğunu gösterir. Kontrol grubunun iki tedavi grubuyla ayrı ayrı
karşılaştırmalarında yeterli kanıt bulunmamıştır.

\begin{mistakebox}
Genel $p=0{,}018$ sonucu üç grubun da birbirinden farklı olduğu anlamına
gelmez. Ayrıca medyanların 4,55, 5,16 ve 5,44 olması Kruskal--Wallis testinin
doğrudan bu üç medyanı sınadığı anlamına gelmez. Test bütün sıra dağılımlarını
kullanır; yalnız konum farkı yorumu dağılım biçimlerinin benzerliğiyle
desteklenmelidir.
\end{mistakebox}

\begin{outputguidebox}
Çıktıyı şu sırayla okuyun: (1) bağımsız grup yapısı, (2) her grubun geçerli
$n$ değeri, medyanı ve kutu grafiği, (3) sıra ortalamaları, (4) genel
Kruskal--Wallis $H$, serbestlik derecesi ve $p$-değeri, (5) epsilon-kare,
(6) yalnız genel sonuç destekliyorsa düzeltilmiş ikili karşılaştırmalar,
(7) karşılaştırma yönü ve dağılım biçimleri. Ham ikili $p$-değerlerini çoklu
test düzeltmesi yapılmadan raporlamayın.
\end{outputguidebox}

\subsection{Örnek bulgu paragrafı}

``Bitki kuru ağırlıklarının sıra dağılımları kontrol ($Md=5{,}16$), tedavi 1
($Md=4{,}55$) ve tedavi 2 ($Md=5{,}44$) gruplarında farklılık göstermiştir,
$H(2)=7{,}99$, $p=0{,}018$, $\varepsilon^2=0{,}222$. Bonferroni düzeltmeli
Dunn karşılaştırmalarında tedavi 2 grubunun sıraları tedavi 1 grubundan daha
yüksektir ($p_{\text{düz.}}=0{,}015$). Kontrol--tedavi 1
($p_{\text{düz.}}=0{,}791$) ve kontrol--tedavi 2
($p_{\text{düz.}}=0{,}273$)
karşılaştırmaları belirgin değildir. Sonuç bütün medyanların birbirinden
farklı olduğu biçiminde yorumlanmamıştır.''

\section{Bölüm sonu çözümlü problemler}

\begin{enumerate}
  \item \textbf{Bağlı sıralar.} $2,4,4,7,9$ değerlerine sıra verin.

  \textbf{Çözüm:} Sıralar $1,2{,}5,2{,}5,4,5$'tir. Eşit iki 4, ikinci ve üçüncü sıraların ortalamasını alır.

  \item \textbf{Mann--Whitney etkisi.} $n_1=10$, $n_2=12$ ve birinci grup yönündeki $U=90$ ise üstünlük olasılığı ve sıra-biserial etkiyi bulun.

  \textbf{Çözüm:} $\hat\theta=U/(n_1n_2)=90/120=0{,}75$'tir. $r_{rb}=2\hat\theta-1=0{,}50$ olur. Birinci gruptan seçilen rastgele bir gözlemin ikinci gruptakinden yüksek olma eğilimi yüzde 75'tir; bağların yarısı hesaba katılır.

  \item \textbf{Kruskal--Wallis etkisi.} $H=12$, $k=3$ ve $N=36$ için epsilon-kareyi hesaplayın.

  \textbf{Çözüm:} $\varepsilon^2=(H-k+1)/(N-k)=(12-3+1)/(36-3)=10/33\approx0{,}303$'tür. Genel sonuç farkın hangi gruplar arasında olduğunu göstermez; düzeltilmiş ikili karşılaştırmalar gerekir.
\end{enumerate}

\bolumkaynaklari{\cite{conover1999,hollander2013,lehmann2006,wasserman2006,efron1993,dobson1983}}

\section*{Hatırla--Uygula--Yansıt}

\begin{enumerate}
\renewcommand{\labelenumi}{\textbf{\Alph{enumi}.}}
  \item \textbf{Hatırla:} Parametrik olmayan yöntemlerin de gerektirdiği dört koşul yazın.
  \item \textbf{Seç:} İki bağımsız sıralı grup, aynı kişilerin iki ölçümü ve üç bağımsız grup için uygun testleri belirleyin.
  \item \textbf{Sırala:} $2,4,4,7,9$ değerlerine bağlı sıraları dikkate alarak sıra verin.
  \item \textbf{Yorumla:} Mann--Whitney için üstünlük olasılığı 0,70 ne anlama gelir?
  \item \textbf{Eleştir:} Mann--Whitney testini her zaman medyan testi saymanın sorununu açıklayın.
  \item \textbf{Ayır:} İşaret testi ile Wilcoxon testinin kullandığı bilgileri ve varsayımları karşılaştırın.
  \item \textbf{Yön:} Fark ``sonra eksi önce'' ve sıra-biserial etki negatifse sonucu yorumlayın.
  \item \textbf{Genel test:} Anlamlı Kruskal--Wallis sonucunun hangi grupların farklı olduğunu neden göstermediğini açıklayın.
  \item \textbf{Etki:} $H=12$, $k=3$, $N=36$ için epsilon-kareyi hesaplayın.
  \item \textbf{Kesinlik:} Kesin ve asimptotik $p$-değerlerinin neden farklılaşabileceğini belirtin.
  \item \textbf{Varsayım:} Bağımsızlık ihlalinin neden sıra testine geçerek çözülmediğini açıklayın.
  \item \textbf{Alternatif:} Ortalama farkı hedefleniyorsa sıra testi dışında üç sağlam yaklaşım önerin.
  \item \textbf{Uygula:} Python verisine aykırı bir değer ekleyerek Welch ve Mann--Whitney sonuçlarını karşılaştırın.
  \item \textbf{Çoklu test:} Kruskal--Wallis sonrası altı ikili karşılaştırmanın neden düzeltme gerektirdiğini açıklayın.
  \item \textbf{Raporla:} Betimsel özet, test istatistiği, yöntem türü, $p$-değeri, etki büyüklüğü ve yorum sınırı içeren paragraf yazın.
\end{enumerate}